\section{Schubert coefficients}

In this section we define Schubert coefficients of matroids and recap what is previously known. For each point $x\in G(r,n)$ we may consider the class of the torus orbit closure $\T$ in $A^\bullet(G(r,n))$. We express this class in the basis of Schubert cycles as
\begin{equation}
    \label{eq:TClass}
    [\T] = \sum_{\lambda} d_\lambda \sigma_\lambda,
\end{equation}
where the sum is over all partitions $\lambda \subseteq r\times(n-r)$ with $|\lambda| = \codim(\T)$ and $d_\lambda \in \Z$. More precisely the coefficient $d_\lambda = \deg([\T] \sigma_{\lambda^c})$ is the number of points in the intersection of $\T$ and the Schubert variety $X_{\lambda^c}(V_\bullet)$ for a generic flag $V_\bullet$. Hence the coefficients are non-negative.

In \cite{Gel} it is shown that $\T$ is isomorphic to the projective toric variety of the matroid polytope $P(M_x)$. Hence many properties of $\T$ can be read off from the matroid $M_x$. For example
\begin{equation}
    \dim(\T) = \dim(P(M_x)) = n - \kappa(M_x)
\end{equation}
and
\begin{equation}
    \deg(\T) = \Vol(P(M_x)).
\end{equation}
Here $\deg(\T)$ denotes the degree of $\T$ embedded in $\PP^{\binom{n}{r}-1}$ and $\Vol(P(M_x))$ is the normalized $n-\kappa(M_x)$ dimensional volume of $P(M_x)$.

As we will see, the class of $\T$ in $A^\bullet(G(r,n))$ only depends on the matroid $M_x$, so we may denote the coefficients in \cref{eq:TClass} by $d_\lambda(M_x)$. The coefficient $d_\lambda(M_x)$ only makes sense if $|\lambda| = r(n-r)-(n-\kappa(M_x))$. For all other cases we define $d_\lambda(M_x) = 0$. This way we can consider $d_\lambda$ as functions from the set of matroids representable over $\CC$ of rank $r$ on $[n]$ to $\Z$. We call the integers $d_\lambda(M_x)$ the Schubert coefficients of the matroid $M_x$.

In \cite{Klyachko}, Klyachko gives a formula for $d_\lambda(U_{r,n})$ in terms of evaluations of Schur polynomials.

\begin{thm}\textnormal{\cite[Theorem 6]{Klyachko}}
    \label{thm:Klyachko}
    Let $\lambda \subseteq r \times (n-r)$ be a partition with $|\lambda| = (r-1)(n-r-1)$. The associated Schubert coefficient of $U_{r,n}$ is
    \begin{equation*}
        d_\lambda(U_{r,n}) = \sum_{i=0}^r (-1)^i\binom{n}{i}s_{\lambda^c}(1^{r-i}).
        \qedhere
    \end{equation*}
\end{thm}

In \cite{Speyer}, Speyer proved the following important theorem describing the Schubert coefficient associated to the hook complement.
\begin{thm}\textnormal{\cite[Theorem 5.1]{Speyer}}
    \label{thm:Speyer}
    For $x \in G(r,n)$
    \begin{equation*}
        d_{h^c}(M_x) = \beta(M_x).
    \end{equation*}
\end{thm}
Notice that this is true even if $M_x$ is disconnected, since in that case both $d_{h^c}(M_x)$ and $\beta(M_x)$ are zero.

One issue with the Schubert coefficients of matroids is that points in the Grassmannian only give rise to matroids representable over $\CC$. This is rectified by Fink and Speyer in \cite{FinkSpeyer} where they define the class of any matroid \(M\) in the \(K\)-theory of the Grassmannian, which can be specialized to a class in the Chow ring. We denote this class \(\Sc(M) \in A^\bullet(G(r,n))\). If \(M\) is represented by a point \(x \in G(r,n)\), then \(\Sc(M) = [\T]\). The specialization from \(K\)-theory to the Chow ring is made explicit by Berget and Fink in \cite{BergetFink}. Note that $\Sc(M) \in A^{r(n-r)-(n-\kappa(M))}(G(r,n))$. From the class $\Sc(M)$ we get Schubert coefficients $d_\lambda(M)$ for arbitrary matroids. We summarize the most relevant takeaways from \cite{FinkSpeyer} and \cite{BergetFink} in a theorem.
\begin{thm}
\label{thm:BergetFink}
    For each partition $\lambda \subseteq r\times (n-r)$ there are functions $d_\lambda \colon S(r,[n]) \to \Z$ coinciding with the coefficients of \cref{eq:TClass} whenever $M = M_x$ for $x\in G(r,n)$. These functions are valuative matroid invariants and whenever $|\lambda| = r(n-r)-(n-1)$ then $d_\lambda$ is additive.
\end{thm}
Notice that since valuative matroid invariants are determined by their values on Schubert matroids, which are representable, it follows from \cref{thm:Speyer} and \cref{thm:BergetFink} that \(d_{h^c}(M) = \beta(M)\) for all matroids \(M\).

As mentioned in \cite{BergetFink} there is another way of defining Schubert coefficients of arbitrary matroids by the theory of tautological classes of matroids from \cite{BEST}. In this setting $d_\lambda(M) = \deg_{X_n}(s_{\lambda^c}(\mathcal{S}_M^\vee))$ where $X_n$ is the $(n-1)$-dimensional permutahedral variety, $\mathcal{S}_M$ is the tautological sub-bundle of $M$ and $s_{\lambda^c}(\mathcal{S}_M^\vee)$ is the evaluation of the Schur polynomial $s_{\lambda^c}$ at the Chern roots of $\mathcal{S}_M^\vee$.

For representable matroids $M_x$ we could conclude that $d_\lambda(M_x)$ is non-negative since it is the number of points in an intersection of varieties. This argument only works for representable matroids, but the non-negativity of $d_\lambda(M)$ for arbitrary matroids $M$ is conjectured in \cite{BergetFink}.
\begin{conj}\textnormal{\cite[Conjecture 9.13]{BergetFink}}
    \label{conj:positivity}
    The Schubert coefficients $d_\lambda(M)$ are non-negative for all matroids $M$.
\end{conj}

